%==================================================================
%  CV - Emmanuel Narro  |  Desarrollador de Software Jr.  (1 PAGINA)
%  Version compacta: sans-serif, acento cobre, niveles, objetivo.
%  Compilar con: pdflatex CV_Emmanuel_Narro_1pagina.tex  (o Overleaf)
%==================================================================
\documentclass[10pt,a4paper]{article}

\usepackage[utf8]{inputenc}
\usepackage[T1]{fontenc}
\usepackage{textcomp}
\usepackage[spanish,es-noshorthands]{babel}
\usepackage[top=1cm,bottom=1cm,left=1.25cm,right=1.25cm]{geometry}
\usepackage{titlesec}
\usepackage{enumitem}
\usepackage{multicol}
\usepackage{xcolor}
\usepackage{hyperref}
\usepackage{helvet}
\renewcommand{\familydefault}{\sfdefault}
\usepackage{microtype}

\definecolor{accent}{HTML}{A85A22}   % cobre elegante
\definecolor{ink}{HTML}{1A2130}
\definecolor{muted}{HTML}{5A6472}

\hypersetup{colorlinks=true, urlcolor=accent, linkcolor=accent}

\titleformat{\section}
  {\color{accent}\bfseries\normalsize}
  {}{0pt}{}[{\color{accent!35}\titlerule[0.8pt]}]
\titlespacing{\section}{0pt}{5pt}{2pt}

\setlist[itemize]{leftmargin=1.1em, itemsep=0.5pt, topsep=1pt, parsep=0pt}
\pagestyle{empty}
\setlength{\parindent}{0pt}
\linespread{0.98}

\newcommand{\entry}[4]{%
  \textbf{#1} \hfill {\color{muted}\small #2}\\[0pt]
  {\color{accent}\small #3} \hfill {\color{muted}\itshape\small #4}\\[1pt]
}

\begin{document}

%==================== ENCABEZADO ====================
\begin{center}
  {\color{ink}\LARGE\bfseries Emmanuel Narro}\\[2pt]
  {\color{accent} Desarrollador de Software Jr. \textbar{} Backend \& Full Stack}\\[3pt]
  {\color{muted}\small
   Guadalajara, Jalisco, México \enspace\textbullet\enspace
   \href{https://wa.me/526182704889}{+52 618 270 4889 (WhatsApp)} \enspace\textbullet\enspace
   \href{mailto:emmanuel.nare.12@gmail.com}{emmanuel.nare.12@gmail.com}\\[1pt]
   \href{https://www.linkedin.com/in/emmanuel-narro-8b4476170}{linkedin.com/in/emmanuel-narro-8b4476170}
   \enspace\textbullet\enspace
   \href{https://github.com/Emmanuel121003}{github.com/Emmanuel121003}
   \enspace\textbullet\enspace
   {\color{accent}\textbf{Portafolio:}}~\href{https://emmanuel121003.github.io/}{emmanuel121003.github.io}}
\end{center}

%==================== PERFIL ====================
\section{Perfil profesional}
Ingeniero en Desarrollo y Gestión de Software con experiencia profesional construyendo y dando soporte a \textbf{sistemas
empresariales en producción} (\textbf{.NET/C\#, Python, SQL Server, APIs REST}). \textbf{Busco integrarme como Desarrollador
Junior Full Stack o Backend}, para aportar en proyectos reales y seguir creciendo en desarrollo backend e integración de sistemas.

%==================== EXPERIENCIA ====================
\section{Experiencia profesional}

\entry{E-GO Desarrollo de Software}{Guadalajara, Jal.}{Desarrollador de Software}{Mayo 2026 -- Actualidad}
\begin{itemize}
  \item \textbf{Scale II} (control de peso industrial, en producción y uso diario en planta): principal contribuidor de la app Android (Java); integración con escritorio .NET/C\#, Raspberry Pi y SQL Server.
  \item \textbf{Gens} (POS + e-commerce): contribuidor del POS en C\#/.NET y del frontend (Vue.js) y servicios web (Node.js).
  \item \textbf{PVM} (punto de venta móvil): plataforma web en PHP y aporte a la app Android.
  \item Disponibilidad y confiabilidad del sistema: despliegue en sitio, incidencias, sync/respaldo por APIs REST y optimización SQL; Scrum, Git y GitHub.
\end{itemize}

\entry{E-GO Desarrollo de Software}{Guadalajara, Jal.}{Becario de Desarrollo de Software}{Enero -- Abril 2026}
\begin{itemize}
  \item Desarrollo del sistema de control de peso e integración entre escritorio, móvil y Raspberry Pi (Python, .NET, SQL Server, Android Studio); pruebas y documentación.
\end{itemize}

\entry{VirtuaMX}{Durango, México}{Becario en Desarrollo Web}{Mayo -- Agosto 2024}
\begin{itemize}
  \item Módulos frontend y backend y APIs REST bajo metodologías ágiles, entregando en tiempo y forma.
\end{itemize}

%==================== PROYECTOS ====================
\section{Proyectos destacados}
\begin{itemize}
  \item \textbf{MedicalRecords} (propio): expedientes médicos en C\#/.NET con SQL Server, pruebas y documentación.
  \item \textbf{Sistema de Asistencia} (PHP/MySQL): roles, registro por QR y huella, modelo relacional y reportes.
  \item \textbf{Sistema Ganadero} (Python/Flask/SQLite): registro de ganado con exportación a Excel y PDF.
  \item \textbf{Scanner-QR} (React Native/Expo) y \textbf{AgroBot} (IoT + IA en equipo: Arduino, Firebase, Vertex AI).
\end{itemize}

%==================== EDUCACIÓN ====================
\section{Educación}
\textbf{Ing. en Desarrollo y Gestión de Software} --- UT de Rodeo \hfill {\color{muted}\small 2024 -- 2026}\\[1pt]
\textbf{TSU en Tecnologías de la Información} (Desarrollo de SW Multiplataforma) --- UT de Rodeo \hfill {\color{muted}\small 2022 -- 2024}

%==================== HABILIDADES ====================
\section{Habilidades técnicas} \textit{\small (A = avanzado, I = intermedio, B = básico)}
\begin{itemize}
  \item \textbf{Backend:} C\#/.NET, PHP, Java \textit{(A)}; Python, ASP.NET Web API, Node.js, APIs REST, Flask \textit{(I)}; Laravel, MVC \textit{(B)}.
  \item \textbf{Frontend:} JavaScript, HTML5, CSS3, Vue.js, TypeScript \textit{(I)}; React, Tailwind \textit{(B)}.
  \item \textbf{Bases de datos:} SQL Server, MySQL \textit{(A)}; SQLite \textit{(I)}; PostgreSQL, MongoDB \textit{(B)}.
  \item \textbf{Móvil / Sistemas:} Java/Android \textit{(A)}; React Native, Expo, Git, GitHub, Postman, Scrum \textit{(I)}; Linux/Raspberry Pi \textit{(B)}.
  \item \textbf{Otros:} Ciberseguridad, ISO 27001, IoT (Arduino/ESP8266), Google Vertex AI \textit{(B)}.
\end{itemize}

%==================== CERTIFICACIONES ====================
\section{Certificaciones}
{\small
\begin{multicols}{2}
\begin{itemize}
  \item Artificial Intelligence Fundamentals --- IBM
  \item Cybersecurity Fundamentals --- IBM
  \item Python Essentials 1 y 2 --- Cisco
  \item SQL y Bases de Datos desde Cero --- Udemy
  \item Introducción a la Ciencia de Datos --- Santander
  \item Internet de las Cosas (IoT) --- Santander
  \item Seguridad digital --- Santander
  \item Microsoft Copilot / IA y Productividad --- Santander
  \item EF SET English --- Nivel B2 (59/100)
\end{itemize}
\end{multicols}}

%==================== IDIOMAS ====================
\section{Idiomas}
\textbf{Español:} nativo. \qquad \textbf{Inglés:} B2 (Intermedio alto), respaldado por EF SET; en desarrollo continuo.

\end{document}
